% paper9_AF_extension.tex
% Extension to the Internal Space: Sigma on A_F
% Author: Sheng-Kai Huang
% This file is a self-contained section (~4 pages)
% to be appended to paper9_spectral_fisher.tex
% between the Discussion and Acknowledgments sections.
%
% Compile: standalone or paste into paper9_spectral_fisher.tex
%
% Requires: same preamble as paper9_spectral_fisher.tex
%   \documentclass[prl,twocolumn,superscriptaddress,longbibliography]{revtex4-2}
%   \usepackage{amsmath,amssymb,amsfonts,bm,hyperref}
%   \newcommand{\kB}{k_{\mathrm{B}}}
%   \newcommand{\Tr}{\mathrm{Tr}}
%   \newcommand{\tr}{\mathrm{tr}}
%   \newcommand{\SvN}{S_{\mathrm{vN}}}
%   \newcommand{\nF}{n_{\mathrm{F}}}

%=======================================================================
\section{Extension to the internal space}
\label{sec:AF}
%=======================================================================

The preceding sections establish $\Sigma = 2\ln Q$ as the
Fisher information of the $a_2$ sector under
\emph{outer} (conformal) deformation of the
metric~[KNOWN, Paper~9 main body].
The matter content of the Standard Model resides in
$a_4$ and is invisible to this conformal probe
($g_{QQ}^{(2)} = 0$).
Here we show how \emph{inner} fluctuations of
the Dirac operator on the finite spectral triple
$(\mathcal{A}_F,\, \mathcal{H}_F,\, D_F)$
give rise to a complementary measure
$\Sigma_F$ that encodes the full gauge and
Yukawa sectors~[NEW].


%-----------------------------------------------------------------------
\subsection{Inner fluctuations and the finite triple}
\label{sec:inner}
%-----------------------------------------------------------------------

In the Chamseddine--Connes
framework~\cite{CC1996,vS2024}, the Standard Model
is encoded in the almost-commutative geometry
$M^4 \times F$, where $F$ is a finite
spectral triple with algebra~[KNOWN]
\begin{equation}\label{eq:AF}
  \mathcal{A}_F
  = \mathbb{C} \oplus \mathbb{H}
    \oplus M_3(\mathbb{C})\,.
\end{equation}
The Hilbert space $\mathcal{H}_F
= \mathbb{C}^{96}$ carries three generations
of fermion representations, and the finite Dirac
operator $D_F$ encodes Yukawa couplings and
the CKM matrix~\cite{vS2024}~[KNOWN].

Gauge fields and the Higgs doublet arise as
\emph{inner fluctuations}~\cite{CC1996,Connes2006}:
\begin{equation}\label{eq:DA}
  D \;\to\; D_A
  = D + A + \varepsilon'\, J A J^{-1}\,,
\end{equation}
where $A = \sum_i a_i [D, b_i]$
$(a_i, b_i \in C^\infty(M) \otimes \mathcal{A}_F)$
and $J$ is the real structure operator of
KO-dimension~$6$~[KNOWN].
The decomposition of $A$ on $M^4 \times F$
yields~\cite{vS2024}~[KNOWN]:
\begin{itemize}
\item a $\mathrm{U}(1)_Y \times \mathrm{SU}(2)_L
\times \mathrm{SU}(3)_C$ gauge field
$A_\mu$ from the manifold directions;
\item a Higgs doublet $H \in (1,2,\tfrac{1}{2})$
from the finite (off-diagonal) directions.
\end{itemize}

The spectral action
$\Tr\, f(D_A^2/\Lambda^2)$, expanded via
the heat kernel, gives the full bosonic
Lagrangian of the Standard Model
coupled to gravity~\cite{CC1996,CC2007,vS2024},
with gauge coupling
unification at $\Lambda$ and the
tree-level prediction
$\sin^2\theta_W = 3/8$~[KNOWN].


%-----------------------------------------------------------------------
\subsection{Definition of $\Sigma_F$}
\label{sec:SigmaF}
%-----------------------------------------------------------------------

We define the internal QRE as
the Umegaki relative entropy between the
Gibbs state of the full (fluctuated) Dirac
operator and the Gibbs state of the free
(unfluctuated) Dirac operator~[NEW]:
\begin{equation}\label{eq:SigmaF_def}
  \boxed{\;
  \Sigma_F
  \;\equiv\;
  D\!\left(\rho_A \,\big\|\, \rho_0\right)
  \;}
\end{equation}
where
\begin{align}
  \rho_A &= e^{-\beta\, D_A^2}\big/\, Z_A\,,
  \label{eq:rhoA}\\
  \rho_0 &= e^{-\beta\, D^{\,2}}\big/\, Z_0\,,
  \label{eq:rho0}
\end{align}
and $\beta = 1/\Lambda^2$ is the spectral
inverse temperature.

By the standard identity for the QRE between
Gibbs states at the same $\beta$~[KNOWN],
\begin{align}\label{eq:SigmaF_expand}
  \Sigma_F
  &= \beta\,\bigl(
    \langle D^2 \rangle_{\!\rho_A}
    - \langle D_A^2 \rangle_{\!\rho_A}
  \bigr)
  + \ln Z_A - \ln Z_0
  \nonumber\\[4pt]
  &= \beta\,\bigl\langle
    D^2 - D_A^2
  \bigr\rangle_{\!\rho_A}
  - \beta\,(F_A - F_0)\,,
\end{align}
where $F = -\beta^{-1}\ln Z$ is the
free energy~[KNOWN].

Substituting
$D_A^2 = D^2 + \{D,A_{\mathrm{in}}\}
+ A_{\mathrm{in}}^2$
(with $A_{\mathrm{in}} \equiv A + JAJ^{-1}$):
\begin{equation}\label{eq:SigmaF_inner}
  \Sigma_F
  = -\beta\,\bigl\langle
    \{D,A_{\mathrm{in}}\}
    + A_{\mathrm{in}}^2
  \bigr\rangle_{\!\rho_A}
  - \beta\,(F_A - F_0)\,.
\end{equation}


%-----------------------------------------------------------------------
\subsection{Variation of $\Sigma_F$ yields
the SM field equations}
\label{sec:variation}
%-----------------------------------------------------------------------

We now show that
$\delta_A \Sigma_F = 0$ reproduces the
Euler--Lagrange equations of the
spectral action~[NEW].

\emph{Proof sketch.}---Consider a
one-parameter family of inner fluctuations
$A(\varepsilon) = \varepsilon\, A^{(1)}$.
The QRE $\Sigma_F(\varepsilon) = D(\rho_\varepsilon \| \rho_0)$
satisfies $\Sigma_F(0) = 0$, and
its first derivative with respect to the
fluctuation parameter is~[KNOWN, standard QIT]
\begin{equation}\label{eq:first_var}
  \frac{d}{d\varepsilon}\,
  D(\rho_\varepsilon \| \rho_0)
  = \Tr\!\left[
    \frac{d\rho_\varepsilon}{d\varepsilon}
    \bigl(\ln\rho_\varepsilon - \ln\rho_0\bigr)
  \right].
\end{equation}

The physical field equations arise from varying
$\Sigma_F$ with respect to the gauge and Higgs
fields at a \emph{fixed} (physical) value
of~$A$:
\begin{equation}\label{eq:variation_principle}
  \delta_A \Sigma_F = 0\,.
\end{equation}

Using Eq.~\eqref{eq:SigmaF_expand} and
the CCSvS identification~\cite{CCSvS2018}
$\SvN(\rho_\beta) = \Tr\, h(\beta D)
= S_{\mathrm{spectral}}$~[KNOWN],
the variation splits into~[NEW]
\begin{align}\label{eq:var_split}
  \delta_A \Sigma_F
  &= \beta\,\delta_A
    \bigl\langle D^2 \bigr\rangle_{\!\rho_A}
  - \delta_A \SvN(\rho_A)
  \nonumber\\
  &= \beta\,\delta_A
    \bigl\langle D^2 \bigr\rangle_{\!\rho_A}
  - \delta_A S_{\mathrm{spectral}}\,.
\end{align}

The first term is a boundary contribution
that vanishes for compact~$M$
(integration by parts on the trace).
Therefore,~[NEW]
\begin{equation}\label{eq:equivalence}
  \delta_A \Sigma_F = 0
  \quad\Longleftrightarrow\quad
  \delta_A S_{\mathrm{spectral}} = 0\,,
\end{equation}
and the SM field equations follow
identically to the standard NCG
derivation~\cite{CC1996,CC2007,vS2024}:
\begin{align}
  \delta_{A_\mu}\, S_{\mathrm{spectral}} = 0
    &\implies \text{Yang--Mills equations}\,,
    \label{eq:YM}\\
  \delta_H\, S_{\mathrm{spectral}} = 0
    &\implies \text{Higgs field equation}\,,
    \label{eq:Higgs}\\
  \delta_g\, S_{\mathrm{spectral}} = 0
    &\implies \text{Einstein equation}\,.
    \label{eq:Einstein}
\end{align}


%-----------------------------------------------------------------------
\subsection{Fisher metric on the gauge parameter space}
\label{sec:gauge_fisher}
%-----------------------------------------------------------------------

The second-order expansion of $\Sigma_F$ in the
inner fluctuation defines the Fisher information
metric on the space of gauge
configurations~[NEW]:
\begin{equation}\label{eq:FisherF}
  \Sigma_F(\varepsilon)
  = \frac{\varepsilon^2}{2}\,
    g_{AA}^{\mathrm{Fisher}}
  + \mathcal{O}(\varepsilon^3)\,,
\end{equation}
where
\begin{equation}\label{eq:gAA}
  g_{AA}^{\mathrm{Fisher}}
  = \frac{d^2}{d\varepsilon^2}\,
    D(\rho_\varepsilon \| \rho_0)
    \Big|_{\varepsilon = 0}
  = \beta^2\,
    \mathrm{Var}(A_{\mathrm{in}}^{(1)})_{\rho_0}\,.
\end{equation}

By the CCSvS identification,
$g_{AA}^{\mathrm{Fisher}}$ equals the
second variation of the spectral
action~[NEW]:
\begin{equation}\label{eq:gAA_spectral}
  g_{AA}^{\mathrm{Fisher}}
  = -\frac{\delta^2\, \SvN}
    {\delta A\, \delta A}
    \bigg|_{A=0}
  + \beta\,
    \frac{\delta^2\,
    \langle D^2 \rangle}
    {\delta A\, \delta A}
    \bigg|_{A=0}\,.
\end{equation}

For the gauge sector,
the dominant contribution comes from the
$a_4$ coefficient of the heat kernel.
Since the spectral action evaluated at $a_4$
contains the Yang--Mills
term~\cite{CC1996,vS2024}~[KNOWN]
\begin{equation}\label{eq:a4_gauge}
  S^{(4)}_{\mathrm{gauge}}
  = \frac{f_0}{2\pi^2}
    \int\! d^4x\,\sqrt{g}\;
    \sum_{i=1}^{3}\,
    \frac{1}{g_i^2}\,
    F_{\mu\nu}^{(i)}\, F^{(i)\mu\nu}\,,
\end{equation}
the Fisher metric decomposes along the
gauge algebra $\mathfrak{g}
= \mathfrak{u}(1) \oplus \mathfrak{su}(2)
\oplus \mathfrak{su}(3)$
as~[NEW]
\begin{equation}\label{eq:gAA_decomp}
  g_{AA}^{\mathrm{Fisher}}
  = g^{\mathrm{U}(1)}
  \oplus g^{\mathrm{SU}(2)}
  \oplus g^{\mathrm{SU}(3)}\,,
\end{equation}
with relative normalizations determined by
the quadratic Casimirs of the representations
appearing in $\mathcal{H}_F$~[KNOWN]:
\begin{equation}\label{eq:coupling_ratio}
  g_3^2 = g_2^2
  = \tfrac{5}{3}\, g_1^2
  \qquad\text{at } \Lambda\,.
\end{equation}


%-----------------------------------------------------------------------
\subsection{Weinberg angle as a Fisher metric ratio}
\label{sec:Weinberg}
%-----------------------------------------------------------------------

The gauge coupling unification
condition~\eqref{eq:coupling_ratio}
has a direct information-geometric
interpretation~[NEW].
Define the electroweak Fisher information
as the sum of the $\mathrm{U}(1)_Y$
and $\mathrm{SU}(2)_L$ contributions:
\begin{equation}\label{eq:g_EW}
  g_{\mathrm{EW}}
  = g^{\mathrm{U}(1)} + g^{\mathrm{SU}(2)}\,.
\end{equation}
The weak mixing angle at the unification
scale is~[NEW interpretation of KNOWN result]
\begin{equation}\label{eq:sin2theta}
  \boxed{\;
  \sin^2\theta_W
  = \frac{g^{\mathrm{U}(1)}}{g_{\mathrm{EW}}}
  = \frac{g_1^2}{g_1^2 + g_2^2}
  = \frac{3/5}{3/5 + 1}
  = \frac{3}{8}
  \;}
\end{equation}
at the GUT scale $\Lambda$.
The value $3/8$ is entirely determined by
the algebra~$\mathcal{A}_F$;
the $\Sigma$ framework does not modify it
but identifies it as the \emph{fraction of
electroweak Fisher information carried by the
hypercharge sector}~[NEW].

Below $\Lambda$, renormalization group running
modifies this to
$\sin^2\theta_W(M_Z) \approx 0.231$,
in agreement with experiment~[KNOWN].


%-----------------------------------------------------------------------
\subsection{Orthogonality of outer and inner sectors}
\label{sec:orthogonality}
%-----------------------------------------------------------------------

The total $\Sigma$ decomposes as~[NEW]
\begin{equation}\label{eq:Sigma_total}
  \Sigma_{\mathrm{total}}
  = \Sigma_{\mathrm{grav}}
  + \Sigma_F\,,
\end{equation}
where $\Sigma_{\mathrm{grav}} = 2\ln Q$
(outer, $a_2$ sector) and
$\Sigma_F = D(\rho_A \| \rho_0)$
(inner, $a_4$ sector).

\emph{Claim}~[NEW]:
The cross-term of the Fisher metric
vanishes,
\begin{equation}\label{eq:cross}
  g_{Q A}
  \;\equiv\;
  \frac{\partial^2\, \Sigma_{\mathrm{total}}}
  {\partial Q\,\partial A}
  \bigg|_{Q=1,\, A=0}
  = 0\,.
\end{equation}

\emph{Proof.}---Under constant conformal
rescaling $g \to Q^2 g$,
the gauge field strength transforms
as~\cite{Wald1984}~[KNOWN]
\begin{equation}
  \hat{F}_{\mu\nu}\, \hat{F}^{\mu\nu}
  = F_{\mu\nu}\, F^{\mu\nu}
\end{equation}
(conformal invariance of Yang--Mills in
$d=4$~[KNOWN]).
Therefore $S_{\mathrm{spectral}}$ evaluated
on the gauge sector is independent
of~$Q$, and all mixed
$Q$-$A$ derivatives vanish.
\hfill$\square$

\emph{Physical meaning.}---Gravity
($\Sigma_{\mathrm{grav}}$) and the
gauge sector ($\Sigma_F$) are
\emph{independent channels of entropy
production} at the Fisher-metric level.
This decoupling is specific to $d=4$:
in $d \neq 4$, the $a_4$ sector is
\emph{not} conformally invariant,
and the mixed Fisher metric would be
nonzero~[NEW].


%-----------------------------------------------------------------------
\subsection{Summary and outlook}
\label{sec:AF_summary}
%-----------------------------------------------------------------------

Table~\ref{tab:bridge} collects the
results of the extension.

\begin{table}[t]
\centering
\caption{Bridge from gravity to the Standard
Model via $\Sigma$ on $M^4 \times F$.
The spectral action serves as the entropy
potential for both sectors.
Outer fluctuations (conformal~$Q$) probe
the $a_2$ sector and yield
$\Sigma_{\mathrm{grav}} = 2\ln Q$;
inner fluctuations (gauge $A$, Higgs $H$)
probe the $a_4$ sector and yield
$\Sigma_F$ whose variation reproduces the
SM field equations.}
\label{tab:bridge}
\begin{ruledtabular}
\begin{tabular}{lcc}
 & \textbf{Outer} (gravity)
 & \textbf{Inner} (SM) \\
\colrule
Deformation
 & $g \to Q^2 g$
 & $D \to D_A$ \\
Parameter
 & $Q$ (lapse)
 & $A_\mu, H$ \\
Heat kernel sector
 & $a_2$
 & $a_4$ \\
Conformal scaling
 & $Q^2$
 & $Q^0$ (invariant) \\
Fisher metric
 & $2/Q^2$
 & SM kinetic terms \\
$\Sigma$
 & $2\ln Q$
 & $D(\rho_A \| \rho_0)$ \\
Field equations
 & Einstein
 & YM + Higgs \\
Cross-coupling
 & \multicolumn{2}{c}{$g_{QA} = 0$ (orthogonal)} \\
\end{tabular}
\end{ruledtabular}
\end{table}

The outer--inner decomposition yields a
unified variational principle~[NEW]:
\begin{equation}\label{eq:master}
  \delta\Sigma_{\mathrm{total}}
  = \delta\Sigma_{\mathrm{grav}}
  + \delta\Sigma_F
  = 0
\end{equation}
recovers the \emph{full} set of
coupled Einstein--Yang-Mills--Higgs
equations.
The orthogonality~\eqref{eq:cross}
ensures that the gravitational and
matter sectors can be extremized
independently at the Fisher (quadratic)
level.

Several questions remain open.

\emph{(i) Lorentzian signature.}---The
CCSvS entropy and the spectral action
are formulated in Euclidean signature.
A Lorentzian extension requires
the Krein spectral triple
framework~\cite{vandenDungen2016,Besnard2017},
which is not yet fully developed
for $M \times F$.

\emph{(ii) The Khronon as dilaton.}---If
the Khronon field $\phi$ is identified with
the Chamseddine--Connes
dilaton~\cite{CC2006}
via $\phi = -\ln Q$, the ghost
condensation mechanism~\cite{BS2024}
would provide a natural scale-breaking
mechanism within the spectral action.
However, the Khronon mass
$\mu_0 \sim H_0/c \sim 10^{-33}\,\mathrm{eV}$
is $\sim 47$ orders of magnitude below
the natural spectral action scale, and
this hierarchy is not explained by the
present framework.

\emph{(iii) Fermion masses.}---The Dirac
operator $D_F$ encodes Yukawa couplings
as its off-diagonal entries, and the spectral
action produces the correct Yukawa
interaction terms~\cite{CC1996,vS2024}.
Within the $\Sigma$ framework, the Yukawa
couplings enter as components of the
internal Fisher metric
$g_{AA}^{\mathrm{Fisher}}$.
Predicting their numerical values
requires additional input beyond the
algebra $\mathcal{A}_F$---potentially
the exceptional Jordan algebra
$J_3(\mathbb{O})$ of
Refs.~\cite{Singh2025,FureyHughes2024},
where the three-eigenvalue structure
could arise from $\Sigma$
extremization~[SPECULATIVE].

\emph{(iv) Number of generations.}---The
dimension selection $(d-2)(d-3) = 2
\Rightarrow d = 4$ has no direct analog
selecting $N_{\mathrm{gen}} = 3$.
The number of generations is an input in
the NCG framework (the rank of $D_F$).
A derivation from $\Sigma$ principles
would require the framework of $J_3(\mathbb{O})$
(uniqueness of the exceptional Jordan
algebra) or $\mathrm{SO}(8)$
triality~\cite{BaezHuerta2010,FureyHughes2024},
and remains an open problem.


%=======================================================================
% Additional references needed for the extension
%=======================================================================

% Add these to the bibliography of paper9_spectral_fisher.tex:
%
% \bibitem{Connes2006}
% A.~Connes,
% ``Noncommutative geometry and the Standard Model with neutrino mixing,''
% J.\ High Energy Phys.\ \textbf{11}, 081 (2006).
%
% \bibitem{CC2007}
% A.~H.~Chamseddine and A.~Connes,
% ``Why the Standard Model,''
% J.\ Geom.\ Phys.\ \textbf{58}, 38 (2008);
% arXiv:0706.3688.
%
% \bibitem{Singh2025}
% T.~P.~Singh,
% ``Fermion mass ratios from the exceptional Jordan algebra,''
% arXiv:2508.10131 (2025).
%
% \bibitem{FureyHughes2024}
% C.~Furey and N.~Hughes,
% ``Three generations and a trio of trialities,''
% Phys.\ Lett.\ B (2025);
% arXiv:2409.17948.
%
% \bibitem{FarnsworthEtAl2026}
% S.~Farnsworth, F.~Finster, N.~Paganini, and T.~P.~Singh,
% ``Causal fermion systems, non-commutative geometry
% and generalized trace dynamics,''
% arXiv:2603.05018 (2026).
